\section*{Chapitre \ref{ch:mx:transaction-cost-analysis} --- Analyse des coûts de transaction}

\begin{solution}{exo:mx:transaction-cost-analysis:1}
Délai $800\times0.02=16$, exécution $800\times0.03=24$, coût d'opportunité $200\times0.10=20$, frais $800\times0.003=2.4$ dollars~: 62.4 dollars, soit 6.24 cents
par action de la cible de 1\,000 actions.
\end{solution}

\begin{solution}{exo:mx:transaction-cost-analysis:2}
Le rapport VWAP indique qu'il a battu la référence d'un cent par action ($100.05-100.06$)~; le rapport d'arrivée indique qu'il a coûté 3 cents par action exécutée, et avec le délai,
le coût d'opportunité et les frais, soit 6.24 cents par action de l'ordre.
\end{solution}

\begin{solution}{exo:mx:transaction-cost-analysis:3}
$-0.5+1.04\times10.0\times\sqrt{0.1}=2.79$ \omterm{def:mx:the-limit-order-book:tick}{pas de cotation} par action.
\end{solution}

\begin{solution}{exo:mx:transaction-cost-analysis:4}
Le desk a confié à l'\omterm{def:mx:benchmark-algorithms:algo}{algorithme d'exécution} urgent les ordres les plus importants (7\,039 actions en moyenne contre 4\,126) et davantage d'ordres passés les jours de forte volatilité~; ces deux facteurs font monter le
coût quel que soit l'algorithme, si bien qu'une partie des 3.36 \omterm{def:mx:the-limit-order-book:tick}{pas de cotation} bruts tient à la difficulté des ordres et non à l'algorithme.
\end{solution}

\begin{solution}{exo:mx:transaction-cost-analysis:5}
Le \omterm{def:mx:transaction-cost-analysis:reversion}{retour post-négociation} se mesure après la dernière exécution. Les dernières exécutions de l'algorithme patient sont souvent des \omterm{def:mx:child-order-placement:risk}{transactions de solde} (clean-up) qui traversent la fourchette en fin de tranche,
après cinq minutes de pression des ordres en attente~; le cours efface alors cette poussée (3.51 \omterm{def:mx:the-limit-order-book:tick}{pas de cotation} au bout de deux minutes). Les dernières tranches de l'algorithme urgent sont les
plus petites de son calendrier concentré sur le début, et ses transactions initiales importantes ont déjà reculé pendant la fenêtre.
\end{solution}

\begin{solution}{exo:mx:transaction-cost-analysis:6}
Lorsque $1.89<\lambda(5.98^2-4.60^2)=14.6\lambda$, soit $\lambda>0.13$ par \omterm{def:mx:the-limit-order-book:tick}{pas de cotation}.
\end{solution}

\begin{solution}{exo:mx:transaction-cost-analysis:7}
2.32 \omterm{def:mx:the-limit-order-book:tick}{pas de cotation}, avec un intervalle à 95~\% allant de 0.79 à 3.85. La participation capte l'essentiel de l'effet de taille, mais pas les journées volatiles que l'algorithme
urgent a aussi reçues plus souvent, si bien qu'une part du biais subsiste (la valeur vraie est 1.89)~; l'intervalle est plus étroit parce que l'estimation utilise moins de
variables de contrôle.
\end{solution}

\begin{solution}{exo:mx:transaction-cost-analysis:8}
Le VWAP sur l'intervalle inclut les transactions propres de l'algorithme et le cours qu'elles ont déplacé~: il dérive avec l'ordre. Les mêmes ordres coûtent 2.42 \omterm{def:mx:the-limit-order-book:tick}{pas de cotation} par action exécutée
contre le prix d'arrivée.
\end{solution}

\begin{solution}{pb:mx:transaction-cost-analysis:1}
\textbf{1.} Le prix d'arrivée, le VWAP sur l'intervalle de l'ordre, le VWAP de la journée, la clôture.
\textbf{2.} Patient 2.42, $-0.29$, 1.90, 2.51~; urgent 3.66, 1.40, 3.39, 4.10 \omterm{def:mx:the-limit-order-book:tick}{pas de cotation} par action exécutée.
\textbf{3.} Le VWAP de sa fenêtre contient ses propres transactions et leur impact.
\textbf{4.} Le prix d'arrivée (ou prix de décision)~: fixé avant que l'ordre n'atteigne le marché, il mesure ainsi l'ensemble du coût de négociation.
\textbf{5.} Demi-fourchette plus $\eta\sigma\sqrt{Q/V}$~; $\eta=1.04$ (erreur type 0.18) avec une fourchette de $-0.5$ \omterm{def:mx:the-limit-order-book:tick}{pas de cotation}.
\textbf{6.} 2.47 prévus contre 2.73 réalisés en moyenne~; une corrélation de 0.25 ordre par ordre.
\textbf{7.} Le délai, la fourchette, l'impact, le timing, le coût d'opportunité et les frais s'additionnent pour former l'implementation shortfall de Perold, car l'impact et le timing se partagent exactement le coût d'exécution.
\textbf{8.} Patient~: délai $-0.35$, fourchette $-0.65$, impact 2.41, timing 0.59, coût d'opportunité 0.10, frais $-0.20$, total 1.90~; urgent~: $-0.15$, 2.43, 1.13,
$-0.07$, 0.16, 0.29, total 3.79.
\textbf{9.} Le cours efface 3.51 \omterm{def:mx:the-limit-order-book:tick}{pas de cotation} (patient) et 1.06 (urgent) en deux minutes~: impact temporaire, pas de \omterm{def:mx:dark-pools-and-blocks:leakage}{fuite d'information}.
\textbf{10.} 53.5~\% des ordres confiés à l'algorithme urgent, dont davantage de gros ordres ou d'ordres passés les jours volatils (7\,039 actions en moyenne contre 4\,126).
\textbf{11.} 3.36 \omterm{def:mx:the-limit-order-book:tick}{pas de cotation}, erreur type 0.76.
\textbf{12.} Coût lié à l'algorithme, volatilité du régime, participation, terme en racine carrée et son interaction avec l'algorithme~; erreurs types
regroupées par journée, parce que le marché de la journée est commun à ses ordres.
\textbf{13.} \emph{Résultat principal}~: urgent moins patient, ajusté de la difficulté, vaut 2.00 \omterm{def:mx:the-limit-order-book:tick}{pas de cotation} par action avec un intervalle à 95~\% allant de 0.32 à 3.68, contre
un écart brut de 3.36 (erreur type 0.76)~; la valeur vraie, obtenue en faisant passer chaque ordre par les deux algorithmes, est 1.89 (0.25).
\textbf{14.} En résumé~: l'ajustement élimine l'essentiel du biais et l'intervalle le couvre. Un ordre réel n'est traité qu'une fois~; son autre issue n'est jamais observée.
\textbf{15.} Patient 0.77 \omterm{def:mx:the-limit-order-book:tick}{pas de cotation} sous la prévision (0.56), urgent 1.50 au-dessus (0.46).
\textbf{16.} Tout confier à l'algorithme patient~: 2.17 \omterm{def:mx:the-limit-order-book:tick}{pas de cotation} par action sur les 50 derniers jours, contre 3.24 pour le routage du desk.
\textbf{17.} Son coût varie moins (écart type 4.60 contre 5.98 \omterm{def:mx:the-limit-order-book:tick}{pas de cotation})~: un gérant qui pénalise suffisamment la variance le préfère.
\textbf{18.} Si les ordres sont assignés au hasard, l'écart brut est sans biais et l'ajustement ne fait que resserrer l'intervalle.
\textbf{19.} Les ordres difficiles expliquent environ 1.4 des 3.4 \omterm{def:mx:the-limit-order-book:tick}{pas de cotation}~; le reste, soit environ 2 \omterm{def:mx:the-limit-order-book:tick}{pas de cotation}, est imputable à l'algorithme, mesuré dans un intervalle que
une année d'ordres de ce type resserrerait.
\textbf{20.} La même référence, fixée avant la négociation, et un ajustement selon la difficulté de chaque ordre.
\end{solution}

\begin{solution}{iq:mx:transaction-cost-analysis:1}
Le VWAP mesure comment le prix de l'ordre se compare à celui du marché pendant l'exécution de l'ordre, y compris sa propre poussée~; le prix d'arrivée mesure ce que le fonds a payé
par rapport au cours au moment où il a décidé. Les 20 points de base correspondent à l'impact et à la dérive~; le gain sur le VWAP montre que l'algorithme a suivi la journée.
\iqlookfor{Quelle référence évolue avec l'ordre.}
\end{solution}

\begin{solution}{iq:mx:transaction-cost-analysis:2}
Fourchette par rapport au prix milieu juste avant chaque exécution~; impact issu d'un modèle pré-négociation (ou du mouvement du marché ajusté par bêta à un indice ou à un secteur), timing en étant le
reliquat~; ou apparier des ordres similaires de sens opposé pour neutraliser le mouvement du marché.
\iqlookfor{Un modèle ou un contrôle du mouvement du marché.}
\end{solution}

\begin{solution}{iq:mx:transaction-cost-analysis:3}
Le nombre d'ordres et le bruit de chacun (erreurs types regroupées par journée ou par titre), la difficulté des ordres de chaque courtier (taille, volatilité,
\omterm{def:mx:the-almgren-chriss-framework:urgency}{urgence}) et la manière dont les ordres ont été assignés~; idéalement, un routage aléatoire.
\iqlookfor{Erreurs types, difficulté, assignation.}
\end{solution}

\begin{solution}{iq:mx:transaction-cost-analysis:4}
Un cours qui efface la poussée de l'ordre une fois celui-ci achevé traduit un impact temporaire~; un cours qui continue sur sa lancée révèle de l'information. Comparer la dérive post-exécution
à une anticipation pré-négociation, par courtier et par plateforme~; une dérive qui commence avant les grosses transactions des \omterm{def:mx:the-almgren-chriss-framework:parent}{ordres enfants} suggère une \omterm{def:mx:dark-pools-and-blocks:leakage}{fuite d'information}.
\iqlookfor{Conventions de signe~; timing par rapport à l'ordre.}
\end{solution}

\begin{solution}{iq:mx:transaction-cost-analysis:5}
\omterm{def:mx:the-almgren-chriss-framework:parent}{Ordre parent}~: instant et prix de décision, prix d'arrivée, instructions, cible, limite, algorithme et paramètres, annulations. \omterm{def:mx:the-almgren-chriss-framework:parent}{Ordre enfant}~: envoi, accusé de réception et exécution
horaires, plateforme, prix, quantité, frais, indicateur de liquidité~; données de marché autour de chacun (prix milieu, fourchette, profondeur, volume) avec des horodatages à la microseconde.
\iqlookfor{Horodatages~; le prix de décision~; données de plateforme au niveau des \omterm{def:mx:the-almgren-chriss-framework:parent}{ordres enfants}.}
\end{solution}

\begin{solution}{iq:mx:transaction-cost-analysis:6}
Demander la composition de l'univers et le modèle de difficulté du fournisseur, vérifier si les ordres du client ont été routés différemment (plus grands, plus urgents,
plus volatils), refaire la comparaison sur les ordres propres du client ajustés avec son propre modèle pré-négociation, et examiner l'intervalle.
\iqlookfor{Composition~; ajustement~; incertitude.}
\end{solution}
